\chapter{Generalized Solutions and Regularity}

\begin{definition}[Divergence-form operator and generalized solution]
\label{gt-def-8-divergence-form-generalized-solution}
\dcref{ch8:8.1--8.4}
\group{ch08-generalized-solutions}
Let $\Omega\subset\mathbb R^n$ be a domain, and let $a^{ij},b^i,c^i,d$ be
measurable. Set
\begin{equation}
Lu=D_i(a^{ij}D_j u+b^iu)+c^iD_i u+du.
\tag{8.1}
\end{equation}
For weakly differentiable $u$ for which the following integral is defined, put
\begin{equation}
\Lop(u,v)=\int_\Omega\bigl((a^{ij}D_j u+b^iu)D_i v
 -(c^iD_i u+du)v\bigr)\,dx.
\tag{8.2}
\end{equation}
The assertions $Lu=0$, $Lu\geq0$, and $Lu\leq0$ hold in the generalized sense
when $\Lop(u,v)=0$, $\Lop(u,v)\leq0$, and $\Lop(u,v)\geq0$, respectively, for
every nonnegative $v\in C_0^1(\Omega)$. If $f^i,g\in L^1_{\mathrm{loc}}(\Omega)$,
then $u$ is a generalized solution of
\begin{equation}
Lu=g+D_if^i
\tag{8.3}
\end{equation}
when
\begin{equation}
\Lop(u,v)=\int_\Omega(f^iD_iv-gv)\,dx
\qquad\forall v\in C_0^1(\Omega).
\tag{8.4}
\end{equation}
\end{definition}

\begin{remark}[Classical and generalized solutions]
\label{gt-rem-8-classical-generalized-equivalence}
\dcref{ch8:8.1--8.4}
\group{ch08-generalized-solutions}
Classical solutions satisfy these generalized identities whenever the integrals
are defined. Conversely, a $C^2$ generalized solution is classical when the
coefficients needed to expand the divergence terms are sufficiently regular.
\end{remark}

\begin{convention}[Ellipticity and coefficient bounds]
\label{gt-conv-8-ellipticity-coefficient-bounds}
\dcref{ch8:8.5--8.6}
\group{ch08-generalized-solutions}
Unless stated otherwise, $L$ is strictly elliptic and has bounded coefficients:
there are $\lambda>0$, $\Lambda>0$, and $\nu\geq0$ such that
\begin{equation}
a^{ij}(x)\xi_i\xi_j\geq\lambda|\xi|^2
\qquad(x\in\Omega,\ \xi\in\mathbb R^n),
\tag{8.5}
\end{equation}
and, for every $x\in\Omega$,
\begin{equation}
\sum_{i,j}|a^{ij}(x)|^2\leq\Lambda^2,
\qquad
\lambda^{-2}\sum_i(|b^i(x)|^2+|c^i(x)|^2)
 +\lambda^{-1}|d(x)|\leq\nu^2.
\tag{8.6}
\end{equation}
\end{convention}

\begin{definition}[Generalized Dirichlet problem]
\label{gt-def-8-generalized-dirichlet-problem}
\dcref{ch8:8.3--8.4}
\group{ch08-generalized-solutions}
For $\varphi\in W^{1,2}(\Omega)$, a generalized solution $u\in W^{1,2}(\Omega)$
of (8.3) has boundary value $\varphi$ when
$u-\varphi\in W_0^{1,2}(\Omega)$.
\end{definition}

\begin{proposition}[Boundedness of the weak form]
\label{gt-prop-8-bounded-weak-form}
\dcref{ch8:8.7}
\group{ch08-generalized-solutions}
Under (8.6),
\begin{equation}
\begin{split}
|\Lop(u,v)|
&\leq\int_\Omega\bigl(|a^{ij}D_juD_iv|+|b^iuD_iv|
 +|c^ivD_iu|+|duv|\bigr)\,dx\\
&\leq C\|u\|_{W^{1,2}(\Omega)}\|v\|_{W^{1,2}(\Omega)}.
\end{split}
\tag{8.7}
\end{equation}
Thus (8.2) and (8.4) extend from $C_0^1(\Omega)$ to
$W_0^{1,2}(\Omega)$, and $u\mapsto\Lop(u,\cdot)$ defines a bounded map
$W^{1,2}(\Omega)\to W_0^{1,2}(\Omega)^*$.
\end{proposition}
\begin{proof}
\sketch Apply Cauchy--Schwarz to each term and use the coefficient bounds (8.6),
then use the density of $C_0^1(\Omega)$ in $W_0^{1,2}(\Omega)$.
\end{proof}

\section*{8.1. The Weak Maximum Principle}

\begin{definition}[Sobolev boundary inequalities]
\label{gt-def-8-sobolev-boundary-inequalities}
\dcref{ch8:8.1}
\group{ch08-generalized-solutions}
For $u,v\in W^{1,2}(\Omega)$, write $u\leq0$ on $\partial\Omega$ when
$u^+\in W_0^{1,2}(\Omega)$, and write $u\leq v$ there when
$(u-v)^+\in W_0^{1,2}(\Omega)$. Define the reversed inequalities by negation and
\[
\sup_{\partial\Omega}u=
\inf\{k\in\mathbb R:u\leq k\text{ on }\partial\Omega\},
\qquad
\inf_{\partial\Omega}u=-\sup_{\partial\Omega}(-u).
\]
For functions continuous near the boundary, this agrees with the pointwise notion.
\end{definition}

\begin{convention}[Weak nonpositivity of the zeroth-order term]
\label{gt-conv-8-weak-zeroth-order-sign}
\dcref{ch8:8.8}
\group{ch08-generalized-solutions}
The generalized sign condition corresponding to $D_ib^i+d\leq0$ is
\begin{equation}
\int_\Omega(dv-b^iD_iv)\,dx\leq0
\qquad\forall v\in C_0^1(\Omega),\quad v\geq0.
\tag{8.8}
\end{equation}
Under (8.6), it also holds for nonnegative $v\in W_0^{1,1}(\Omega)$.
\end{convention}

\begin{theorem}[Weak maximum principle]
\label{gt-thm-8-1-weak-maximum-principle}
\source{gilbarg-trudinger:page-0191}
\source{gilbarg-trudinger:page-0192}
\dcref{ch8:8.1}
\group{ch08-generalized-solutions}
Assume (8.5), (8.6), and (8.8). If $u\in W^{1,2}(\Omega)$ satisfies
$Lu\geq0$ (respectively, $Lu\leq0$), then
\begin{equation}
\sup_\Omega u\leq\sup_{\partial\Omega}u^+,
\qquad
\left(\text{respectively, }\inf_\Omega u\geq\inf_{\partial\Omega}u^-\right).
\tag{8.9}
\end{equation}
\end{theorem}
\begin{proof}
\sketch It suffices to treat the subsolution case. Put
$l=\sup_{\partial\Omega}u^+$ and test with $v=(u-k)^+$ for
$l\leq k<\sup_\Omega u$. Condition (8.8), applied to $uv$, removes the
zeroth-order contribution, while ellipticity and (8.6) give
\begin{equation}
\tag{8.10}
\int_\Omega a^{ij}D_j uD_i v\,dx
\leq2\lambda\int_\Omega v|Du|\,dx.
\end{equation}
This yields an estimate of
$\|Dv\|_2$ by $\|v\|_{L^2(\supp Dv)}$. Sobolev and H\"older inequalities then
force $|\supp Dv|$ to have a positive lower bound independent of $k$. Letting
$k\uparrow\sup_\Omega u$ contradicts the vanishing of $Du$ on the top level set.
Apply the result to $-u$ for the supersolution statement.
\end{proof}

\begin{corollary}[Uniqueness in $W_0^{1,2}$]
\label{gt-cor-8-2-uniqueness-zero}
\source{gilbarg-trudinger:page-0192}
\dcref{ch8:8.2}
\group{ch08-generalized-solutions}
If $u\in W_0^{1,2}(\Omega)$ and $Lu=0$ in $\Omega$, then $u=0$ in $\Omega$.
\end{corollary}
\begin{proof}
\sketch Apply Theorem 8.1 to $u$ and $-u$.
\end{proof}

\section*{8.2. Solvability of the Dirichlet Problem}

\begin{theorem}[Solvability of the generalized Dirichlet problem]
\label{gt-thm-8-3-solvability-generalized-dirichlet}
\source{gilbarg-trudinger:page-0193}
\dcref{ch8:8.3}
\group{ch08-generalized-solutions}
Assume (8.5), (8.6), and (8.8). For every $\varphi\in W^{1,2}(\Omega)$ and
$g,f^i\in L^2(\Omega)$, the problem
\[
Lu=g+D_if^i\quad\text{in }\Omega,
\qquad u=\varphi\quad\text{on }\partial\Omega
\]
has a unique generalized solution.
\end{theorem}
\begin{proof}
\sketch Set $w=u-\varphi$ and absorb $L\varphi$ into
$\hat g=g-c^iD_i\varphi-d\varphi$ and
$\hat f^i=f^i-a^{ij}D_j\varphi-b^i\varphi$, reducing to zero boundary data.
For sufficiently large $\sigma_0$, Lemma 8.4 makes the form of
$L_{\sigma_0}=L-\sigma_0$ coercive. The equation is then a compact perturbation
of the invertible operator $L_{\sigma_0}$ by Lemma 8.5. The Fredholm alternative
and Corollary 8.2 give existence and uniqueness.
\end{proof}

\begin{lemma}[Coercivity estimate]
\label{gt-lem-8-4-coercivity-estimate}
\source{gilbarg-trudinger:page-0193}
\dcref{ch8:8.4}
\group{ch08-generalized-solutions}
If $L$ satisfies (8.5) and (8.6), then
\begin{equation}
\Qop(u,u)\geq\frac{\lambda}{2}\int_\Omega|Du|^2\,dx
-\lambda\nu^2\int_\Omega u^2\,dx.
\tag{8.11}
\end{equation}
\end{lemma}
\begin{proof}
\sketch Expand $\Qop(u,u)$, use ellipticity on the principal term, and absorb
$(b^i-c^i)uD_iu$ by Young's inequality and (8.6).
\end{proof}

\begin{lemma}[Compact embedding]
\label{gt-lem-8-5-compact-embedding}
\source{gilbarg-trudinger:page-0194}
\dcref{ch8:8.5}
\group{ch08-generalized-solutions}
For $\Hcal=W_0^{1,2}(\Omega)$, the map $I:\Hcal\to\Hcal^*$ defined by
\begin{equation}
Iu(v)=\int_\Omega uv\,dx
\tag{8.12}
\end{equation}
is compact.
\end{lemma}
\begin{proof}
\sketch Factor $I$ through the compact embedding
$W_0^{1,2}(\Omega)\hookrightarrow L^2(\Omega)$ and the continuous map
$L^2(\Omega)\to W_0^{1,2}(\Omega)^*$ induced by the $L^2$ pairing.
\end{proof}

Define $L_\sigma=L-\sigma$. For a coercive shift $\sigma_0$, the zero-boundary
equation is equivalently
\begin{equation}
u+\sigma_0L_{\sigma_0}^{-1}Iu=L_{\sigma_0}^{-1}F.
\tag{8.13}
\end{equation}
The formal adjoint is
\begin{equation}
L^*u=D_i(a^{ij}D_ju-c^iu)-b^iD_iu+du.
\tag{8.14}
\end{equation}

\begin{theorem}[Fredholm alternative]
\label{gt-thm-8-6-fredholm-alternative}
\dcref{ch8:8.6}
\group{ch08-generalized-solutions}
Assume (8.5) and (8.6). There is a countable discrete set
$\Sigma\subset\mathbb R$ such that, if $\sigma\notin\Sigma$, the Dirichlet
problems for $L_\sigma$ and $L_\sigma^*$ are uniquely solvable for arbitrary
$g,f^i\in L^2(\Omega)$ and $\varphi\in W^{1,2}(\Omega)$. If
$\sigma\in\Sigma$, the homogeneous kernels of $L_\sigma$ and $L_\sigma^*$
with zero boundary values have positive finite dimension, and
$L_\sigma u=g+D_if^i$, $u=\varphi$ on $\partial\Omega$, is solvable exactly when
\begin{equation}
\int_\Omega\bigl((g-c^iD_i\varphi-d\varphi+\sigma\varphi)v
 -(f^i-a^{ij}D_j\varphi-b^i\varphi)D_iv\bigr)\,dx=0
\tag{8.15}
\end{equation}
for every $v\in W_0^{1,2}(\Omega)$ satisfying $L_\sigma^*v=0$. If (8.8)
also holds, then $\Sigma\subset(-\infty,0)$.
\end{theorem}
\begin{proof}
\sketch For a coercive shift $L_{\sigma_0}$, rewrite the equation as
$u+(\sigma_0-\sigma)L_{\sigma_0}^{-1}Iu=L_{\sigma_0}^{-1}F$. Lemma 8.5
makes the perturbation compact, and its Hilbert-space adjoint is
$(\sigma_0-\sigma)(L_{\sigma_0}^*)^{-1}I$. Apply the Fredholm alternative,
including its orthogonality criterion. The maximum principle excludes
nonnegative exceptional shifts.
\end{proof}

\begin{corollary}[Resolvent estimate]
\label{gt-cor-8-7-resolvent-estimate}
\dcref{ch8:8.7}
\group{ch08-generalized-solutions}
If $\sigma\notin\Sigma$ and $u\in W^{1,2}(\Omega)$ solves
$L_\sigma u=g+D_if^i$ with $u=\varphi$ on $\partial\Omega$, then
\begin{equation}
\|u\|_{W^{1,2}(\Omega)}
\leq C\left(\|g\|_2+\sum_i\|f^i\|_2
+\|\varphi\|_{W^{1,2}(\Omega)}\right),
\tag{8.16}
\end{equation}
where $C$ depends only on $L$, $\sigma$, and $\Omega$.
\end{corollary}
\begin{proof}
\sketch The Green operator $G_\sigma=L_\sigma^{-1}:\Hcal^*\to\Hcal$ is bounded.
\end{proof}

\begin{remark}[Adjoint sign condition for solvability]
\label{gt-rem-8-adjoint-sign-solvability}
\dcref{ch8:8.3}
\group{ch08-generalized-solutions}
Theorem 8.3 also remains valid when condition (8.8) is replaced by its adjoint
version obtained by replacing $b^i$ with $-c^i$.
\end{remark}

\section*{8.3. Differentiability of Weak Solutions}

\begin{theorem}[Interior $W^{2,2}$ regularity]
\label{gt-thm-8-8-interior-w22}
\dcref{ch8:8.8}
\group{ch08-generalized-solutions}
Let $u\in W^{1,2}(\Omega)$ solve $Lu=f$. Assume strict ellipticity,
$a^{ij},b^i\in C^{0,1}(\Omega)$ uniformly, $c^i,d\in L^\infty(\Omega)$, and
$f\in L^2(\Omega)$. Then $u\in W^{2,2}(\Omega')$ for every
$\Omega'\Subset\Omega$, with
\begin{equation}
\|u\|_{W^{2,2}(\Omega')}
\leq C\bigl(\|u\|_{W^{1,2}(\Omega)}+\|f\|_{L^2(\Omega)}\bigr),
\tag{8.17}
\end{equation}
where $C=C(n,\lambda,K,d')$ and $d'=\dist(\Omega',\partial\Omega)$.
Moreover,
\begin{equation}
Lu=a^{ij}D_{ij}u+(D_ja^{ij}+b^i+c^i)D_iu+(D_ib^i+d)u=f
\tag{8.18}
\end{equation}
almost everywhere.
\end{theorem}
\begin{proof}
\sketch Move the lower-order terms to an $L^2$ function in the weak identity.
Thus
\begin{equation}
\tag{8.19}
\int_\Omega a^{ij}D_j uD_i v\,dx=\int_\Omega gv\,dx,
\qquad v\in C_0^1(\Omega),
\end{equation}
where
\begin{equation}
\tag{8.20}
g=(b^i+c^i)D_i u+(D_i b^i+d)u-f.
\end{equation}
Test its difference quotient in direction $e_k$ against
$\eta^2\Delta_k^hu$. Ellipticity, the Lipschitz coefficient bounds, and Young's
inequality give an $h$-uniform local $L^2$ bound for $D\Delta_k^hu$.
The difference-quotient criterion yields $Du\in W^{1,2}_{\mathrm{loc}}$, and
integration by parts gives (8.18) almost everywhere.
\end{proof}

\begin{theorem}[Strong Dirichlet solvability]
\label{gt-thm-8-9-strong-dirichlet-solvability}
\dcref{ch8:8.9}
\group{ch08-generalized-solutions}
Let
\begin{equation}
Lu=a^{ij}D_{ij}u+b^iD_iu+cu
\tag{8.21}
\end{equation}
be strictly elliptic, with
$a^{ij}\in C^{0,1}(\Omega)$, $b^i,c\in L^\infty(\Omega)$, and $c\leq0$.
For every $f\in L^2(\Omega)$ and $\varphi\in W^{1,2}(\Omega)$ there is a unique
$u\in W^{1,2}(\Omega)\cap W^{2,2}_{\mathrm{loc}}(\Omega)$ such that
$Lu=f$ and $u-\varphi\in W_0^{1,2}(\Omega)$.
\end{theorem}
\begin{proof}
\sketch Rewrite the nondivergence operator in divergence form, apply Theorem 8.3,
and invoke Theorem 8.8 for interior second derivatives.
\end{proof}

\begin{remark}[Continuous principal coefficients]
\label{gt-rem-8-continuous-principal-coefficients}
\dcref{ch8:8.9}
\group{ch08-generalized-solutions}
For sufficiently smooth $\partial\Omega$ and $\varphi\in W^{2,2}(\Omega)$,
the same conclusion holds with only $a^{ij}\in C^0(\overline\Omega)$; this is
the later global result cited as Theorem 9.15 in the baseline.
\end{remark}

\begin{example}[Failure with discontinuous principal coefficients]
\label{gt-ex-8-discontinuous-principal-nonuniqueness}
\dcref{ch8:8.22}
\group{ch08-generalized-solutions}
Let $B=B_1(0)$, $0<\lambda<1$, and
$b=-1+(n-1)/(1-\lambda)$. If $n>2(2-\lambda)>2$, then
\begin{equation}
\Delta u+b\frac{x_ix_j}{|x|^2}D_{ij}u=0
\tag{8.22}
\end{equation}
has the distinct solutions $u_1(x)=1$ and $u_2(x)=|x|^\lambda$ in
$W^{2,2}(B)$, with equal boundary values. Thus Theorem 8.9 need not remain
unique for merely $L^\infty$ principal coefficients.
\end{example}

\begin{theorem}[Higher interior differentiability]
\label{gt-thm-8-10-higher-interior-regularity}
\dcref{ch8:8.10}
\group{ch08-generalized-solutions}
Let $k\geq1$, and let $u\in W^{1,2}(\Omega)$ solve $Lu=f$. Assume strict
ellipticity, $a^{ij},b^i\in C^{k,1}(\overline\Omega)$,
$c^i,d\in C^{k-1,1}(\overline\Omega)$, and $f\in W^{k,2}(\Omega)$.
Then $u\in W^{k+2,2}(\Omega')$ for every $\Omega'\Subset\Omega$, and
\begin{equation}
\|u\|_{W^{k+2,2}(\Omega')}
\leq C\bigl(\|u\|_{W^{1,2}(\Omega)}+\|f\|_{W^{k,2}(\Omega)}\bigr),
\tag{8.24}
\end{equation}
where $C=C(n,\lambda,K,d',k)$ with the stated coefficient norms absorbed in $K$.
\end{theorem}
\begin{proof}
\sketch Differentiating the weak identity gives
\begin{equation}
\tag{8.23}
\int_\Omega a^{ij}D_{jk}uD_i v\,dx
=\int_\Omega D_k\widehat g\,v\,dx
\qquad(v\in C_0^1(\Omega)).
\end{equation}
Thus each $D_ku$ solves an equation of the
same type with $L^2_{\mathrm{loc}}$ data; iterate Theorem 8.8 and its estimate.
\end{proof}

\begin{corollary}[Interior smoothness]
\label{gt-cor-8-11-interior-smoothness}
\dcref{ch8:8.11}
\group{ch08-generalized-solutions}
If the coefficients of a strictly elliptic $L$ and the datum $f$ are in
$C^\infty(\Omega)$, then every weak solution $u\in W^{1,2}(\Omega)$ of $Lu=f$
belongs to $C^\infty(\Omega)$.
\end{corollary}
\begin{proof}
\sketch Combine Theorem 8.10 with Sobolev embedding at successively higher orders.
\end{proof}

\section*{8.4. Global Regularity}

\begin{theorem}[Global $W^{2,2}$ regularity]
\label{gt-thm-8-12-global-w22}
\dcref{ch8:8.12}
\group{ch08-generalized-solutions}
In addition to the hypotheses of Theorem 8.8, suppose
$\partial\Omega\in C^2$ and $u-\varphi\in W_0^{1,2}(\Omega)$ for some
$\varphi\in W^{2,2}(\Omega)$. Then $u\in W^{2,2}(\Omega)$ and
\begin{equation}
\|u\|_{W^{2,2}(\Omega)}\leq C\bigl(\|u\|_{L^2(\Omega)}
 +\|f\|_{L^2(\Omega)}+\|\varphi\|_{W^{2,2}(\Omega)}\bigr),
\tag{8.25}
\end{equation}
where $C=C(n,\lambda,K,\partial\Omega)$.
\end{theorem}
\begin{proof}
\sketch Subtract $\varphi$ and use the coercivity estimate to control the
$W^{1,2}$ norm:
\begin{equation}
\tag{8.26}
\|u\|_{W^{1,2}(\Omega)}
\leq C(\|u\|_2+\|f\|_2).
\end{equation}
Flatten the boundary by $C^2$ charts. Tangential difference
quotients give every second derivative except the double normal derivative,
which (8.18) controls directly. Combine finitely many boundary charts with the
interior estimate.
\end{proof}

\begin{theorem}[Higher global regularity]
\label{gt-thm-8-13-higher-global-regularity}
\dcref{ch8:8.13}
\group{ch08-generalized-solutions}
Under the hypotheses of Theorem 8.10, also assume
$\partial\Omega\in C^{k+2}$ and
$u-\varphi\in W_0^{1,2}(\Omega)$ for some
$\varphi\in W^{k+2,2}(\Omega)$. Then $u\in W^{k+2,2}(\Omega)$ and
\begin{equation}
\|u\|_{W^{k+2,2}(\Omega)}\leq C\bigl(\|u\|_{L^2(\Omega)}
 +\|f\|_{W^{k,2}(\Omega)}+\|\varphi\|_{W^{k+2,2}(\Omega)}\bigr).
\tag{8.27}
\end{equation}
Here $C=C(n,\lambda,K,k,\partial\Omega)$. If the coefficients, $f$, $\varphi$,
and $\partial\Omega$ are smooth, then $u\in C^\infty(\overline\Omega)$.
\end{theorem}
\begin{proof}
\sketch Iterate the tangential differentiation and normal recovery argument from
Theorem 8.12, then apply Sobolev embedding in the smooth case.
\end{proof}

\begin{theorem}[Smooth classical Dirichlet problem]
\label{gt-thm-8-14-smooth-classical-dirichlet}
\dcref{ch8:8.14}
\group{ch08-generalized-solutions}
Let $Lu=a^{ij}D_{ij}u+b^iD_iu+cu$ be strictly elliptic with
$C^\infty(\overline\Omega)$ coefficients and $c\leq0$. If
$\partial\Omega\in C^\infty$, then for arbitrary
$f,\varphi\in C^\infty(\overline\Omega)$ the Dirichlet problem
$Lu=f$, $u=\varphi$ has a unique solution in $C^\infty(\overline\Omega)$.
\end{theorem}
\begin{proof}
\sketch Combine Theorem 8.3 with Theorem 8.13.
\end{proof}

\section*{8.5. Global Boundedness of Weak Solutions}

For
\begin{equation}
D_iA^i(x,u,Du)+B(x,u,Du)=0.
\tag{8.28}
\end{equation}
Here
\begin{equation}
A^i(x,z,p)=a^{ij}p_j+b^iz-f^i,
\qquad B(x,z,p)=c^ip_i+dz-g.
\tag{8.29}
\end{equation}
a weak subsolution (respectively, supersolution or solution) is a weakly
differentiable $u$ such that
\begin{equation}
\int_\Omega(D_ivA^i(x,u,Du)-vB(x,u,Du))\,dx\leq0
\quad(\text{respectively, }\geq0\text{ or }=0)
\tag{8.30}
\end{equation}
for all nonnegative $v\in C_0^1(\Omega)$.

The basic structural estimates are
\begin{equation}
p_iA^i(x,z,p)\geq\frac\lambda2|p|^2
-\frac1\lambda(|bz|^2+|f|^2),
\qquad
|B(x,z,p)|\leq|c|\,|p|+|dz|+|g|.
\tag{8.31}
\end{equation}

For $k>0$, set
\begin{equation}
\bar z=|z|+k,
\qquad
\bar b=\lambda^{-2}(|b|^2+|c|^2+k^{-2}|f|^2)
 +\lambda^{-1}(|d|+k^{-1}|g|).
\tag{8.32}
\end{equation}
The ellipticity and coefficient bounds imply, for $0<\varepsilon<1$,
\begin{equation}
p_iA^i(x,z,p)\geq\frac\lambda2(|p|^2-2\bar b\bar z^2),
\qquad
|zB(x,z,p)|\leq\frac\lambda2
\left(\varepsilon|p|^2+\frac{\bar b}{\varepsilon}\bar z^2\right).
\tag{8.33}
\end{equation}

\begin{theorem}[Global boundedness]
\label{gt-thm-8-15-global-boundedness}
\dcref{ch8:8.15}
\group{ch08-generalized-solutions}
Assume (8.5), (8.6), $f^i\in L^q(\Omega)$, and
$g\in L^{q/2}(\Omega)$ for some $q>n$. If
$u\in W^{1,2}(\Omega)$ is a subsolution with $u\leq0$ on $\partial\Omega$, then
\begin{equation}
\sup_\Omega u\leq C(\|u^+\|_2+k),
\qquad
k=\lambda^{-1}(\|f\|_q+\|g\|_{q/2}),
\tag{8.34}
\end{equation}
where $C=C(n,\nu,q,|\Omega|)$. The corresponding supersolution estimate follows
by replacing $u$ with $-u$.
\end{theorem}
\begin{proof}
\sketch Put $w=u^++k$ and, for a truncated power $H$, test with
\begin{equation}
\tag{8.35}
v=G(w)=\int_k^w|H'(s)|^2\,ds.
\end{equation}
The structural inequalities and Sobolev's inequality give
\begin{equation}
\tag{8.36}
\|H(w)\|_{2\widehat n/(\widehat n-2)}
\leq C\|wH'(w)\|_{2q/(q-2)}.
\end{equation}
After removing the truncation,
\begin{equation}
\tag{8.37}
\|w\|_{\beta\chi q^*}\leq(C\beta)^{1/\beta}\|w\|_{\beta q^*},
\qquad q^*=\frac{2q}{q-2},\quad \chi>1.
\end{equation}
Iterate with $\beta=\chi^m$, pass to $L^\infty$, and interpolate down to $L^2$.
\end{proof}

\begin{remark}[Nonzero constant boundary bound]
\label{gt-rem-8-constant-boundary-bound}
\dcref{ch8:8.38}
\group{ch08-generalized-solutions}
If $u\leq l$ on $\partial\Omega$, the same argument applied to $u-l$ gives
\begin{equation}
\sup_\Omega u\leq C(\|u\|_2+\bar k+|l|),
\qquad
\bar k=k+\lambda^{-1}|l|(\|b\|_q+\|d\|_{q/2}).
\tag{8.38}
\end{equation}
The supersolution statement is obtained by negation, and solutions satisfy the
corresponding estimate for $|u|$.
\end{remark}

\begin{theorem}[Global a priori maximum bound]
\label{gt-thm-8-16-global-apriori-bound}
\dcref{ch8:8.16}
\group{ch08-generalized-solutions}
Under the hypotheses of Theorem 8.15 and condition (8.8), every subsolution
(respectively, supersolution) satisfies
\begin{equation}
\sup_\Omega u\leq\sup_{\partial\Omega}u^++Ck
\qquad
\left(\text{respectively, }-\inf_\Omega u\leq
\sup_{\partial\Omega}u^-+Ck\right),
\tag{8.39}
\end{equation}
where $k$ and $C=C(n,\nu,q,|\Omega|)$ are as above.
\end{theorem}
\begin{proof}
\sketch Subtract the boundary supremum. The weak inequality becomes
\begin{equation}
\tag{8.40}
\int_\Omega\bigl(a^{ij}D_j uD_i v-(b^i+c^i)vD_i u\bigr)\,dx
\leq\int_\Omega(f^iD_i v-gv)\,dx
\end{equation}
for admissible nonnegative $v$. Set $M=\sup_\Omega u^+$ and test with
$u^+/(M+k-u^+)$. This bounds the
$W^{1,2}$ norm of $w=\log((M+k)/(M+k-u^+))$. A second transformed test shows
\begin{equation}
\tag{8.41}
\|w\|_2\leq C(n,\nu,|\Omega|).
\end{equation}
It also gives the weak inequality
\begin{equation}
\tag{8.42}
\int_\Omega\bigl(a^{ij}D_jwD_i\eta-(b^i+c^i)\eta D_iw\bigr)\,dx
\leq\int_\Omega(\widehat g\eta+\widehat f^iD_i\eta)\,dx,
\end{equation}
where
$\widehat g=|g|/k+|f|^2/(2\lambda k^2)$ and
$\widehat f^i=f^i/(M+k-u^+)$. Thus $w$ is a subsolution with uniformly
controlled data.
Theorem 8.15 bounds
$\sup w$, hence $(M+k)/k$, uniformly. Apply the same argument to $-u$.
\end{proof}

\begin{remark}[Adjoint sign maximum bound]
\label{gt-rem-8-adjoint-sign-maximum-bound}
\dcref{ch8:8.16}
\group{ch08-generalized-solutions}
Theorem 8.16 also holds with the adjoint sign condition obtained by replacing
$b^i$ with $-c^i$ in (8.8).
\end{remark}

\section*{8.6. Local Properties of Weak Solutions}

Besides (8.31)--(8.33), the local arguments use
\begin{equation}
|A(x,z,p)|\leq|a|\,|p|+|bz|+|f|.
\tag{8.43}
\end{equation}
After dividing by $\lambda/2$, these inequalities may be written
\begin{equation}
\begin{gathered}
|A(x,z,p)|\leq|a|\,|p|+2\bar b^{1/2}\bar z,
\qquad p\cdot A(x,z,p)\geq|p|^2-2\bar b\bar z^2,\\
|\bar zB(x,z,p)|\leq\varepsilon|p|^2
+\varepsilon^{-1}\bar b\bar z^2
\qquad(0<\varepsilon\leq1).
\end{gathered}
\tag{8.44}
\end{equation}

For a radius $R>0$ and $q>n$, write
\begin{equation}
k(R)=\lambda^{-1}\bigl(R^\delta\|f\|_q+R^{2\delta}\|g\|_{q/2}\bigr),
\qquad \delta=1-\frac nq.
\tag{8.45}
\end{equation}

\begin{theorem}[Local boundedness]
\label{gt-thm-8-17-local-boundedness}
\dcref{ch8:8.17}
\group{ch08-generalized-solutions}
Assume (8.5), (8.6), $f^i\in L^q(\Omega)$, and
$g\in L^{q/2}(\Omega)$ for some $q>n$. If $u\in W^{1,2}(\Omega)$ is a
subsolution, then for every $B_{2R}(y)\subset\Omega$ and $p>1$,
\begin{equation}
\sup_{B_R(y)}u^+\leq C\left(R^{-n/p}\|u^+\|_{L^p(B_{2R}(y))}+k(R)\right),
\tag{8.46}
\end{equation}
where $C=C(n,\Lambda/\lambda,\nu R,q,p)$. The supersolution version follows by
replacing $u$ with $-u$.
\end{theorem}
\begin{proof}
\sketch After scaling to $R=1$, use the test
$\eta^2(u^++k)^\beta$ with $\beta>0$. The structural estimates give a Caccioppoli
inequality for $(u^++k)^{(\beta+1)/2}$. Sobolev iteration on nested balls raises
the integrability exponent to infinity, producing (8.46); rescaling restores $R$.
\end{proof}

\begin{theorem}[Weak Harnack inequality]
\label{gt-thm-8-18-weak-harnack}
\dcref{ch8:8.18}
\group{ch08-generalized-solutions}
Under the coefficient and data hypotheses of Theorem 8.17, let $u$ be a
supersolution that is nonnegative in $B_{4R}(y)\subset\Omega$. For
$1\leq p<n/(n-2)$,
\begin{equation}
R^{-n/p}\|u\|_{L^p(B_{2R}(y))}
\leq C\left(\inf_{B_R(y)}u+k(R)\right),
\tag{8.47}
\end{equation}
where $C=C(n,\Lambda/\lambda,\nu R,q,p)$.
\end{theorem}
\begin{proof}
\sketch Scale first to $R=1$, put $\bar u=u+k$, and use, for a cutoff
$\eta$ and $\beta\ne0$,
\begin{equation}
\tag{8.48}
v=\eta^2\bar u^\beta,
\end{equation}
whose derivative is
\begin{equation}
\tag{8.49}
Dv=2\eta D\eta\,\bar u^\beta
+\beta\eta^2\bar u^{\beta-1}Du.
\end{equation}
Substitution in the weak inequality gives
\begin{equation}
\tag{8.50}
\beta\int\eta^2\bar u^{\beta-1}Du\cdot A
+2\int\eta D\eta\cdot A\,\bar u^\beta
-\int\eta^2\bar u^\beta B\ \begin{cases}\leq0,&u\text{ subsolution},\\
\geq0,&u\text{ supersolution}.
\end{cases}
\end{equation}
For $0<\varepsilon\leq1$, the structural terms satisfy
\begin{equation}
\tag{8.51}
|\eta D\eta\cdot A\,\bar u^\beta|
\leq\frac\varepsilon2\eta^2\bar u^{\beta-1}|Du|^2
+\left(1+\frac{|a|^2}{2\varepsilon}\right)|D\eta|^2\bar u^{\beta+1}
+b\eta^2\bar u^{\beta+1}.
\end{equation}
Taking $\varepsilon=\min(1,|\beta|/4)$ yields
\begin{equation}
\tag{8.52}
\int\eta^2\bar u^{\beta-1}|Du|^2
\leq C(|\beta|)\int
\bigl(b\eta^2+(1+|a|^2)|D\eta|^2\bigr)\bar u^{\beta+1}.
\end{equation}
With $\gamma=\beta+1$ and
$w=\bar u^{\gamma/2}$ for $\gamma\ne0$, while $w=\log\bar u$ for
$\gamma=0$, this becomes
\begin{equation}
\tag{8.53}
\int|\eta Dw|^2\leq
\begin{cases}
C(|\beta|)\gamma^2\int
(b\eta^2+(1+|a|^2)|D\eta|^2)w^2,&\gamma\ne0,\\
C\int(b\eta^2+(1+|a|^2)|D\eta|^2),&\gamma=0.
\end{cases}
\end{equation}
Sobolev, H\"older, and interpolation inequalities give
\begin{equation}
\tag{8.54}
\|\eta w\|_{2\widehat n/(\widehat n-2)}
\leq C(1+|\gamma|)^{\sigma+1}
\|(\eta+|D\eta|)w\|_2,
\end{equation}
where $\sigma=\widehat n/(q-\widehat n)$. For cutoffs between
$B_{r_1}$ and $B_{r_2}$,
\begin{equation}
\tag{8.55}
\|w\|_{2\chi;B_{r_1}}
\leq\frac{C(1+|\gamma|)^{\sigma+1}}{r_2-r_1}
\|w\|_{2;B_{r_2}},
\qquad\chi=\frac{\widehat n}{\widehat n-2}.
\end{equation}
Introduce
\begin{equation}
\tag{8.56}
\Phi(p,r)=\left(\int_{B_r}|\bar u|^p\,dx\right)^{1/p}.
\end{equation}
Then (8.55) is
\begin{equation}
\tag{8.57}
\begin{aligned}
\Phi(\chi\gamma,r_1)&\leq
\left(\frac{C(1+|\gamma|)^{\sigma+1}}{r_2-r_1}\right)^{2/|\gamma|}
\Phi(\gamma,r_2),&&\gamma>0,\\
\Phi(\gamma,r_2)&\leq
\left(\frac{C(1+|\gamma|)^{\sigma+1}}{r_2-r_1}\right)^{2/|\gamma|}
\Phi(\chi\gamma,r_1),&&\gamma<0.
\end{aligned}
\end{equation}
Positive iteration gives
\begin{equation}
\tag{8.58}
\sup_{B_1}\bar u\leq C\|\bar u\|_{L^p(B_2)},
\end{equation}
while negative iteration gives, for $0<p_0<p<\chi$,
\begin{equation}
\tag{8.59}
\Phi(p,2)\leq C\Phi(p_0,3),
\qquad\Phi(-p_0,3)\leq C\Phi(-\infty,1).
\end{equation}
It remains to bridge the moments:
\begin{equation}
\tag{8.60}
\Phi(p_0,3)\leq C\Phi(-p_0,3).
\end{equation}
For $w=\log\bar u$, the second line of (8.53) implies on every
$B_{2r}\subset B_4$
\begin{equation}
\tag{8.61}
\int_{B_r}|Dw|\,dx\leq Cr^{n-1}.
\end{equation}
John--Nirenberg gives (8.60). Rescale and let the auxiliary positive $k$
decrease to $k(R)$.
\end{proof}

\section*{8.7. The Strong Maximum Principle}

\begin{theorem}[Strong maximum principle]
\label{gt-thm-8-19-strong-maximum-principle}
\dcref{ch8:8.19}
\group{ch08-generalized-solutions}
Assume (8.5), (8.6), and (8.8), and let $u\in W^{1,2}(\Omega)$ satisfy
$Lu\geq0$. If some ball $B\Subset\Omega$ satisfies
\begin{equation}
\sup_Bu=\sup_\Omega u\geq0,
\tag{8.62}
\end{equation}
then $u$ is constant in $\Omega$. If this constant is nonzero, equality holds in
(8.8).
\end{theorem}
\begin{proof}
\sketch On a concentric ball compactly contained in $\Omega$, apply Theorem 8.18
with $p=1$ to the nonnegative supersolution $M-u$, where $M=\sup_\Omega u$.
Its infimum vanishes, so $u=M$ on a larger ball. Propagate this equality through
the connected domain. Substitution of a nonzero constant gives equality in (8.8).
\end{proof}

\section*{8.8. The Harnack Inequality}

\begin{theorem}[Harnack inequality]
\label{gt-thm-8-20-harnack}
\dcref{ch8:8.20}
\group{ch08-generalized-solutions}
Assume (8.5) and (8.6). If $u\in W^{1,2}(\Omega)$ is nonnegative and
$Lu=0$, then for every $B_{4R}(y)\subset\Omega$,
\begin{equation}
\sup_{B_R(y)}u\leq C\inf_{B_R(y)}u,
\tag{8.63}
\end{equation}
where $C=C(n,\Lambda/\lambda,\nu R)$. One may take
$C\leq C_0^{\Lambda/\lambda+\nu R}$ with $C_0=C_0(n)$.
\end{theorem}
\begin{proof}
\sketch Apply Theorem 8.17 to control the supremum by a positive $L^p$ norm, and
Theorem 8.18 to control that norm by the infimum.
\end{proof}

\begin{corollary}[Harnack inequality on compact subdomains]
\label{gt-cor-8-21-harnack-subdomain}
\dcref{ch8:8.21}
\group{ch08-generalized-solutions}
Under the hypotheses of Theorem 8.20, every $\Omega'\Subset\Omega$ satisfies
\begin{equation}
\sup_{\Omega'}u\leq C\inf_{\Omega'}u,
\tag{8.64}
\end{equation}
where $C=C(n,\Lambda/\lambda,\nu,\Omega',\Omega)$.
\end{corollary}
\begin{proof}
\sketch Join points of $\Omega'$ by a finite chain of interior balls and iterate
Theorem 8.20.
\end{proof}

\section*{8.9. H\"older Continuity}

\begin{theorem}[Interior oscillation estimate]
\label{gt-thm-8-22-interior-holder}
\dcref{ch8:8.22}
\group{ch08-generalized-solutions}
Assume (8.5), (8.6), $f^i\in L^q(\Omega)$, and
$g\in L^{q/2}(\Omega)$ for some $q>n$. If $u\in W^{1,2}(\Omega)$ solves
(8.3), then $u$ is locally H\"older continuous. For
$B_0=B_{R_0}(y)\subset\Omega$ and $R\leq R_0$,
\begin{equation}
\osc_{B_R(y)}u\leq CR^\alpha
\left(R_0^{-\alpha}\sup_{B_0}|u|+k\right),
\qquad k=\lambda^{-1}(\|f\|_q+\|g\|_{q/2}),
\tag{8.65}
\end{equation}
where $C=C(n,\Lambda/\lambda,\nu,q,R_0)$ and
$\alpha=\alpha(n,\Lambda/\lambda,\nu R_0,q)>0$.
\end{theorem}
\begin{proof}
\sketch Apply the weak Harnack inequality to
$M_{4R}-u$ and $u-m_{4R}$. Adding the resulting estimates gives
$\omega(R)\leq\gamma\omega(4R)+k(R)$ for the oscillation $\omega$ and a fixed
$0<\gamma<1$. Lemma 8.23 converts this decay recurrence into (8.65).
\end{proof}

\begin{lemma}[Iteration of an oscillation inequality]
\label{gt-lem-8-23-oscillation-iteration}
\dcref{ch8:8.23}
\group{ch08-generalized-solutions}
Let $\omega$ and $\sigma$ be nondecreasing on $(0,R_0]$, and suppose
\begin{equation}
\omega(\tau R)\leq\gamma\omega(R)+\sigma(R)
\tag{8.66}
\end{equation}
for $0<\gamma,\tau<1$. For every $0<\mu<1$ and $R\leq R_0$,
\begin{equation}
\omega(R)\leq C\left(
\left(\frac R{R_0}\right)^\alpha\omega(R_0)
+\sigma(R^\mu R_0^{1-\mu})\right),
\tag{8.67}
\end{equation}
where $C=C(\gamma,\tau)$ and $\alpha=\alpha(\gamma,\tau,\mu)>0$.
\end{lemma}
\begin{proof}
\sketch Iterate (8.66) from an intermediate radius $R_1$ to obtain geometric
decay plus $\sigma(R_1)/(1-\gamma)$. Choose
$R_1=R^\mu R_0^{1-\mu}$ and compare $R$ with adjacent powers of $\tau R_1$.
\end{proof}

\begin{theorem}[Interior H\"older norm estimate]
\label{gt-thm-8-24-interior-holder-norm}
\dcref{ch8:8.24}
\group{ch08-generalized-solutions}
Under the hypotheses of Theorem 8.22, for every $\Omega'\Subset\Omega$,
\begin{equation}
\|u\|_{C^\alpha(\overline{\Omega'})}
\leq C(\|u\|_{L^2(\Omega)}+k),
\tag{8.68}
\end{equation}
where $d'=\dist(\Omega',\partial\Omega)$,
$C=C(n,\Lambda/\lambda,\nu,q,d')$,
$\alpha=\alpha(n,\Lambda/\lambda,\nu d')>0$, and
$k=\lambda^{-1}(\|f\|_q+\|g\|_{q/2})$.
\end{theorem}
\begin{proof}
\sketch Use Theorem 8.22 with radius comparable to $d'$ and Theorem 8.17 to
bound the local supremum by the global $L^2$ norm and the data.
\end{proof}

\begin{remark}[Dependence of local-estimate constants]
\label{gt-rem-8-local-constant-dependence}
\dcref{ch8:8.46--8.68}
\group{ch08-generalized-solutions}
The constants in (8.46), (8.47), and (8.63) are nondecreasing in $\nu R$;
the constant in (8.65) is nondecreasing in $R_0$; and the exponents in
(8.65) and (8.68) are nonincreasing in $\nu R_0$ and $\nu d'$, respectively.
When $\nu=0$, these constants and exponents are independent of the displayed
radius and hence of the corresponding domains.
\end{remark}

\section*{8.10. Local Estimates at the Boundary}

\begin{definition}[Sobolev inequality on a boundary portion]
\label{gt-def-8-boundary-portion-inequality}
\dcref{ch8:8.25--8.26}
\group{ch08-generalized-solutions}
For $T\subset\overline\Omega$ and $u\in W^{1,2}(\Omega)$, write $u\leq0$ on
$T$ when $u^+$ is a $W^{1,2}(\Omega)$ limit of functions in
$C_0^1(\overline\Omega\setminus T)$. The other inequalities are defined as in
Section 8.1. This agrees with pointwise ordering for continuous functions and
with the earlier definition when $T=\partial\Omega$.
\end{definition}

\begin{theorem}[Boundary local boundedness]
\label{gt-thm-8-25-boundary-local-boundedness}
\dcref{ch8:8.25}
\group{ch08-generalized-solutions}
Under the hypotheses of Theorem 8.17, let $u$ be a subsolution. For
$y\in\mathbb R^n$, $R>0$, and $p>1$, set
\[
M=\sup_{\partial\Omega\cap B_{2R}(y)}u^+,
\qquad
u_M^+(x)=
\begin{cases}
\max\{u(x),M\},&x\in\Omega,\\
M,&x\notin\Omega.
\end{cases}
\]
Then
\begin{equation}
\sup_{B_R(y)}u_M^+
\leq C\left(R^{-n/p}\|u_M^+\|_{L^p(B_{2R}(y))}+k(R)\right),
\tag{8.69}
\end{equation}
where $C=C(n,\Lambda/\lambda,\nu R,q,p)$.
\end{theorem}
\begin{proof}
\sketch Extend the boundary truncation by the constant $M$ and use the positive
power tests from Theorem 8.17 after subtracting their boundary value. Their
support lies where the structural inequalities apply, so the same iteration gives
(8.69).
\end{proof}

\begin{theorem}[Boundary weak Harnack inequality]
\label{gt-thm-8-26-boundary-weak-harnack}
\dcref{ch8:8.26}
\group{ch08-generalized-solutions}
Under the hypotheses of Theorem 8.18, let $u$ be a supersolution nonnegative in
$\Omega\cap B_{4R}(y)$. Put
\[
m=\inf_{\partial\Omega\cap B_{4R}(y)}u,
\qquad
u_m^-(x)=
\begin{cases}
\min\{u(x),m\},&x\in\Omega,\\
m,&x\notin\Omega.
\end{cases}
\]
For $1\leq p<n/(n-2)$,
\begin{equation}
R^{-n/p}\|u_m^-\|_{L^p(B_{2R}(y))}
\leq C\left(\inf_{B_R(y)}u_m^-+k(R)\right),
\tag{8.70}
\end{equation}
where $C=C(n,\Lambda/\lambda,\nu R,q,p)$.
\end{theorem}
\begin{proof}
\sketch Use negative-power tests for $u_m^-+k$, subtracting the constant boundary
value so the tests vanish in the Sobolev sense. Together with the positive
boundary truncation, the shared test is
\begin{equation}
\tag{8.71}
v=\eta^2
\begin{cases}
\bar u^\beta-(M+k)^\beta,&\beta>0,\\
\bar u^\beta-(m+k)^\beta,&\beta<0.
\end{cases}
\end{equation}
The proof of Theorem 8.18 then
applies unchanged.
\end{proof}

\begin{definition}[Exterior cone condition]
\label{gt-def-8-exterior-cone-condition}
\dcref{ch8:8.27}
\group{ch08-generalized-solutions}
The domain $\Omega$ satisfies an exterior cone condition at
$x_0\in\partial\Omega$ when a finite right circular cone $V_{x_0}$ with vertex
$x_0$ satisfies $\overline\Omega\cap V_{x_0}=\{x_0\}$. It satisfies a uniform
exterior cone condition on $T\subset\partial\Omega$ when these cones may be
chosen congruent to one fixed cone.
\end{definition}

\begin{theorem}[Boundary oscillation estimate]
\label{gt-thm-8-27-boundary-oscillation}
\dcref{ch8:8.27}
\group{ch08-generalized-solutions}
Under the hypotheses of Theorem 8.22, suppose $\Omega$ has an exterior cone at
$x_0\in\partial\Omega$. For $B_0=B_{R_0}(x_0)$ and $0<R\leq R_0$,
\begin{equation}
\osc_{\Omega\cap B_R(x_0)}u
\leq C\left(R^\alpha
\left(R_0^{-\alpha}\sup_{\Omega\cap B_0}|u|+k\right)
+\sigma(\sqrt{RR_0})\right),
\tag{8.72}
\end{equation}
where $\sigma(r)=\osc_{\partial\Omega\cap B_r(x_0)}u$ and $C,\alpha>0$
depend only on $n,\Lambda/\lambda,\nu,q,R_0$, and the cone $V_{x_0}$ as in
the baseline statement.
\end{theorem}
\begin{proof}
\sketch Apply Theorem 8.26 to $M_{4R}-u$ and $u-m_{4R}$. The exterior cone gives
a scale-independent lower bound for the relative measure outside $\Omega$,
producing an oscillation recurrence with an additional boundary-oscillation
term. Lemma 8.23 yields (8.72).
\end{proof}

\begin{corollary}[Uniform continuity up to the boundary]
\label{gt-cor-8-28-uniform-boundary-continuity}
\dcref{ch8:8.28}
\group{ch08-generalized-solutions}
In addition to the hypotheses of Theorem 8.27, suppose the exterior cone
condition holds at every boundary point and
\[
\osc_{\partial\Omega\cap B_R(x_0)}u\longrightarrow0
\qquad(R\downarrow0)
\]
for every $x_0\in\partial\Omega$. Then $u$ is uniformly continuous in $\Omega$.
\end{corollary}
\begin{proof}
\sketch Combine the boundary modulus from Theorem 8.27 with the interior
H\"older estimate of Theorem 8.22 and compactness of $\partial\Omega$.
\end{proof}

\begin{theorem}[H\"older estimate on a boundary portion]
\label{gt-thm-8-29-boundary-holder}
\dcref{ch8:8.29}
\group{ch08-generalized-solutions}
Assume the hypotheses of Theorem 8.22 and a uniform exterior cone condition on
$T\subset\partial\Omega$. If constants $K,\alpha_0>0$ satisfy
\[
\osc_{\partial\Omega\cap B_R(x_0)}u\leq KR^{\alpha_0}
\qquad(x_0\in T,\ R>0),
\]
then $u\in C^\alpha(\Omega\cup T)$ for some $\alpha>0$, and every
$\Omega'\Subset\Omega\cup T$ satisfies
\begin{equation}
\|u\|_{C^\alpha(\Omega')}
\leq C(\sup_\Omega|u|+K+k),
\tag{8.73}
\end{equation}
with $d'=\dist(\Omega',\partial\Omega\setminus T)$ and the dependence of
$C,\alpha$ on $n,\Lambda/\lambda,\nu,V,q,\alpha_0,d'$ stated in the baseline.
For $\Omega'=\Omega$, replace $d'$ by $\diam\Omega$.
\end{theorem}
\begin{proof}
\sketch Use Theorem 8.22 on the ball whose radius is the distance to the boundary.
Theorem 8.27 controls its oscillation by the global supremum, the data, and the
boundary H\"older modulus. In particular,
\begin{equation}
\tag{8.74}
\frac{|u(x)-u(y)|}{|x-y|^\alpha}
\leq C(\sup_\Omega|u|+k+K)
\end{equation}
whenever $x$ lies in the ball centred at $y$ of radius
$\dist(y,\partial\Omega)$. Applying the boundary estimate again when two points
are farther apart than that radius produces a uniform seminorm.
\end{proof}

\begin{remark}[Boundary oscillation from Sobolev trace data]
\label{gt-rem-8-boundary-oscillation-trace}
\dcref{ch8:8.29}
\group{ch08-generalized-solutions}
If $u-v\in W_0^{1,2}(\Omega)$ and $v\in C^0(\overline\Omega)$, then the boundary
oscillation of $u$ tends to zero at every boundary point. If
$v\in C^{\alpha_0}(\overline\Omega)$, it satisfies a uniform bound of order
$R^{\alpha_0}$.
\end{remark}

\begin{theorem}[Continuous Dirichlet data under an exterior cone condition]
\label{gt-thm-8-30-continuous-dirichlet-cone}
\dcref{ch8:8.30}
\group{ch08-generalized-solutions}
Assume (8.5), (8.6), (8.8), $f^i\in L^q(\Omega)$, and
$g\in L^{q/2}(\Omega)$ for some $q>n$. If $\Omega$ has an exterior cone at every
boundary point, then for every $\varphi\in C^0(\partial\Omega)$ there is a unique
\[
u\in W^{1,2}_{\mathrm{loc}}(\Omega)\cap C^0(\overline\Omega)
\]
satisfying $Lu=g+D_if^i$ in $\Omega$ and $u=\varphi$ on $\partial\Omega$.
\end{theorem}
\begin{proof}
\sketch Approximate $\varphi$ uniformly by traces of $C^1(\overline\Omega)$
functions and solve the corresponding generalized Dirichlet problems by
Theorem 8.3. The maximum principle makes the solutions uniformly Cauchy;
Corollary 8.28 gives boundary continuity, and the local energy estimate gives
$W^{1,2}_{\mathrm{loc}}$ convergence. Pass to the weak equation. Local use of
Theorem 8.1 gives uniqueness.
\end{proof}

\begin{definition}[Regular boundary point]
\label{gt-def-8-regular-boundary-point}
\dcref{ch8:8.75}
\group{ch08-generalized-solutions}
A point $x_0\in\partial\Omega$ is regular for $L$ if every solution constructed
from continuous boundary data and admissible $f^i,g$ converges to its prescribed
boundary value as $x\to x_0$. The local density condition
\begin{equation}
\liminf_{R\downarrow0}\frac{|B_R(x_0)\setminus\Omega|}{R^n}>0
\tag{8.75}
\end{equation}
is sufficient by the proof of Theorem 8.27.
\end{definition}

\begin{remark}[Operator independence of regular points]
\label{gt-rem-8-regular-points-laplacian}
\dcref{ch8:8.31}
\group{ch08-generalized-solutions}
Under the chapter's ellipticity and coefficient hypotheses, the regular points
for $L$ coincide with the regular points for the Laplacian.
\end{remark}

\begin{theorem}[Continuous Dirichlet data under the Wiener condition]
\label{gt-thm-8-31-continuous-dirichlet-wiener}
\dcref{ch8:8.31}
\group{ch08-generalized-solutions}
Assume (8.5), (8.6), (8.8), $f^i\in L^q(\Omega)$, and
$g\in L^{q/2}(\Omega)$ for some $q>n$. If the Wiener condition (2.37) holds at
every point of $\partial\Omega$, then for each $\varphi\in C^0(\partial\Omega)$
there is a unique
$u\in W^{1,2}_{\mathrm{loc}}(\Omega)\cap C^0(\overline\Omega)$ satisfying
$Lu=g+D_if^i$ and $u=\varphi$ on $\partial\Omega$.
\end{theorem}
\begin{proof}
\sketch The negative-power estimate underlying Theorem 8.18 also controls
the gradient through
\begin{equation}
\tag{8.76}
R^{1-n}\int_{B_{2R}}|Du|
\leq\left(R^{-n}\int_{B_{2R}}\bar u^p\right)^{1/2}
\left(R^{2-n}\int_{B_{2R}}\bar u^{-p}|Du|^2\right)^{1/2}
\leq C(\inf_{B_R}u+k(R)).
\end{equation}
For the boundary truncation this becomes
\begin{equation}
\tag{8.77}
R^{1-n}\int_{\Omega\cap B_{2R}}|Du_m|
\leq C(\inf_{\Omega\cap B_R}u_m+k(R)).
\end{equation}
With a cutoff $\eta$ and $w=\eta u_m^-$, an energy test gives
\begin{equation}
\tag{8.78}
R^{2-n}\int_{B_{2R}}|Dw|^2
\leq C(m+k)(\inf_{B_R}u_m^-+k(R)).
\end{equation}
The variational characterization is
\begin{equation}
\tag{8.79}
\operatorname{cap}(B_R\setminus\Omega)
=\inf_{v\in\mathcal K}\int|Dv|^2,
\qquad
\mathcal K=\{v\in C_0^1(\mathbb R^n):v=1\text{ on }B_R\setminus\Omega\}.
\end{equation}
Combining (8.77)--(8.79) gives
\begin{equation}
mR^{2-n}\operatorname{cap}(B_R\setminus\Omega)
\leq C\inf_{B_R}(u_m^-+k(R)).
\tag{8.80}
\end{equation}
Writing $\chi(R)=R^{2-n}\operatorname{cap}(B_R\setminus\Omega)$ yields
\begin{equation}
\osc_{\Omega\cap B_R}u
\leq\left(1-\frac{\chi(R)}C\right)\osc_{\Omega\cap B_{4R}}u
+\frac{\chi(R)}C\osc_{\partial\Omega\cap B_{4R}}u+k(R).
\tag{8.81}
\end{equation}
The Wiener divergence condition forces this iteration to inherit the continuous
boundary modulus. Existence and uniqueness then follow by the approximation
argument of Theorem 8.30.
\end{proof}

\begin{remark}[Integrable lower-order coefficients]
\label{gt-rem-8-integrable-lower-order-coefficients}
\dcref{ch8:8.17--8.31}
\group{ch08-generalized-solutions}
The conclusions of Sections 8.6--8.10 remain valid if the boundedness part of
(8.6) for $b,c,d$ is replaced by
$b,c\in L^q(\Omega)$ and $d\in L^{q/2}(\Omega)$ with $q>n$. With
\[
\mathfrak b=\lambda^{-2}(|b|^2+|c|^2)+\lambda^{-1}|d|,
\qquad \nu^2=\|\mathfrak b\|_{q/2},
\]
replace $\nu R$ in Theorems 8.17--8.29 by $\nu R^\delta$,
$\delta=1-n/q$.
\end{remark}

\section*{8.11. H\"older Estimates for the First Derivatives}

The constant-coefficient starting point is
\begin{equation}
\Delta u=g+D_if^i.
\tag{8.82}
\end{equation}
For $g,f\in L^\infty(\Omega)$ and $f\in C^\alpha(\Omega)$, its Newtonian
potential
\[
w(x)=\int_\Omega\Gamma(x-y)g(y)\,dy
+\int_\Omega D_i\Gamma(x-y)f^i(y)\,dy
\]
is a weak solution. Consider now
\begin{equation}
Lu=g+D_if^i.
\tag{8.83}
\end{equation}
Freeze the principal coefficients at $x_0$:
\begin{equation}
\begin{split}
a^{ij}(x_0)D_{ij}u
&=D_i\bigl((a^{ij}(x_0)-a^{ij}(x))D_ju-b^iu+f^i\bigr)\\
&\quad-c^iD_iu-du+g=G+D_iF^i.
\end{split}
\tag{8.84}
\end{equation}
Assume in this section that $0<\alpha<1$, $L$ satisfies (8.5),
$a^{ij},b^i\in C^\alpha(\overline\Omega)$,
$c^i,d,g\in L^\infty(\Omega)$, $f\in C^\alpha(\overline\Omega)$, and
\begin{equation}
\max_{i,j}\{|a^{ij},b^i|_{0,\alpha;\Omega},|c^i,d|_{0;\Omega}\}\leq K.
\tag{8.85}
\end{equation}

\begin{theorem}[Interior gradient H\"older estimate]
\label{gt-thm-8-32-interior-gradient-holder}
\dcref{ch8:8.32}
\group{ch08-generalized-solutions}
If $u\in C^{1,\alpha}(\Omega)$ is a weak solution of
$Lu=g+D_if^i$ in a bounded domain, then for every $\Omega'\Subset\Omega$,
\begin{equation}
|u|_{1,\alpha;\Omega'}
\leq C(|u|_{0;\Omega}+|g|_{0;\Omega}+|f|_{0,\alpha;\Omega}),
\tag{8.86}
\end{equation}
where $C=C(n,\lambda,K,\dist(\Omega',\partial\Omega))$.
\end{theorem}
\begin{proof}
\sketch Apply the constant-coefficient interior estimate to (8.84). On sufficiently
small balls, the H\"older oscillation of $a^{ij}$ absorbs the term containing $Du$;
cover and interpolate to obtain (8.86).
\end{proof}

\begin{theorem}[Global gradient H\"older estimate]
\label{gt-thm-8-33-global-gradient-holder}
\dcref{ch8:8.33}
\group{ch08-generalized-solutions}
If $\Omega$ is $C^{1,\alpha}$, $u\in C^{1,\alpha}(\overline\Omega)$ solves
$Lu=g+D_if^i$, and $u=\varphi$ on $\partial\Omega$ for
$\varphi\in C^{1,\alpha}(\overline\Omega)$, then
\begin{equation}
|u|_{1,\alpha}
\leq C(|u|_0+|\varphi|_{1,\alpha}+|g|_0+|f|_{0,\alpha}),
\tag{8.87}
\end{equation}
where $C=C(n,\lambda,K,\partial\Omega)$.
\end{theorem}
\begin{proof}
\sketch Flatten the boundary by $C^{1,\alpha}$ charts, apply the corresponding
constant-coefficient boundary estimate to (8.84), absorb the coefficient
oscillation locally, and combine a finite boundary cover with Theorem 8.32.
\end{proof}

\begin{theorem}[$C^{1,\alpha}$ Dirichlet solvability]
\label{gt-thm-8-34-c1alpha-dirichlet}
\dcref{ch8:8.34}
\group{ch08-generalized-solutions}
Let $\Omega$ be $C^{1,\alpha}$ and suppose $L$ satisfies (8.5), (8.8), and
(8.85). For $g\in L^\infty(\Omega)$,
$f^i\in C^\alpha(\overline\Omega)$, and
$\varphi\in C^{1,\alpha}(\overline\Omega)$, the generalized Dirichlet problem
\[
Lu=g+D_if^i\quad\text{in }\Omega,
\qquad u=\varphi\quad\text{on }\partial\Omega
\]
has a unique solution in $C^{1,\alpha}(\overline\Omega)$.
\end{theorem}
\begin{proof}
\sketch Approximate the coefficients, data, boundary values, and domain by smooth
objects while preserving ellipticity, (8.8), and uniform H\"older bounds. Solve the
smooth problems. Theorem 8.16 controls the uniform norm and Theorem 8.33 gives
a uniform $C^{1,\alpha}$ bound. Compactness and passage through the weak identity
give a limit solution; Theorem 8.1 gives uniqueness even in the associated
$W^{1,2}$ Dirichlet class.
\end{proof}

\begin{corollary}[$W^{1,2}$ solutions satisfy the global gradient estimate]
\label{gt-cor-8-35-w12-global-gradient-estimate}
\dcref{ch8:8.35}
\group{ch08-generalized-solutions}
Under the hypotheses of Theorem 8.33, if $u\in W^{1,2}(\Omega)$ solves (8.83)
and $u-\varphi\in W_0^{1,2}(\Omega)$, then $u\in C^{1,\alpha}(\overline\Omega)$
and satisfies (8.87).
\end{corollary}
\begin{proof}
\sketch Theorem 8.16 first bounds $u$. For a sufficiently large $\sigma>0$,
Theorem 8.34 solves the shifted equation
$(L-\sigma)v=g+D_if^i-\sigma u$ in $C^{1,\alpha}$. Uniqueness for the shifted
problem gives $v=u$, and Theorem 8.33 applies.
\end{proof}

\begin{corollary}[Local and partial-boundary gradient regularity]
\label{gt-cor-8-36-partial-boundary-gradient-regularity}
\dcref{ch8:8.36}
\group{ch08-generalized-solutions}
Let $T$ be a possibly empty $C^{1,\alpha}$ portion of $\partial\Omega$. If
$u\in W^{1,2}(\Omega)$ solves (8.83) and $u=0$ on $T$ in the Sobolev sense,
then $u\in C^{1,\alpha}(\Omega\cup T)$ and, for every
$\Omega'\Subset\Omega\cup T$,
\begin{equation}
|u|_{1,\alpha;\Omega'}
\leq C(|u|_{0;\Omega}+|g|_{0;\Omega}+|f|_{0,\alpha;\Omega}),
\tag{8.91}
\end{equation}
where $C=C(n,\lambda,K,d')$ and
$d'=\dist(\Omega',\partial\Omega\setminus T)$. If $u=\varphi$ on $T$ for
$\varphi\in C^{1,\alpha}(\overline\Omega)$, add
$|\varphi|_{1,\alpha;\Omega}$ to the right-hand side.
\end{corollary}
\begin{proof}
\sketch Localize the approximation argument for Corollary 8.35, using interior
charts away from $T$ and flattened boundary charts along $T$. For nonzero data,
apply the zero-boundary statement to $u-\varphi$.
\end{proof}

\begin{remark}[Integrable scalar datum]
\label{gt-rem-8-integrable-scalar-datum}
\dcref{ch8:8.32--8.36}
\group{ch08-generalized-solutions}
The conclusions of this section remain valid with
$g\in L^p(\Omega)$, $p=n/(1-\alpha)$, when the $L^p$ norm of $g$ replaces its
supremum norm in the estimates.
\end{remark}

\section*{8.12. The Eigenvalue Problem}

Assume here that $L$ is self-adjoint:
\[
Lu=D_i(a^{ij}D_ju+b^iu)-b^iD_iu+cu,
\qquad a^{ij}=a^{ji}.
\]
On $H=W_0^{1,2}(\Omega)$ define
\[
\Lop(u,u)=\int_\Omega(a^{ij}D_iuD_ju+2b^iuD_iu+cu^2)\,dx,
\qquad
J(u)=\frac{\Lop(u,u)}{\|u\|_2^2}\quad(u\neq0).
\]
The first variational value and its Euler equation are
\begin{equation}
\sigma_1=\inf_{u\in H\setminus\{0\}}J(u),
\tag{8.92}
\end{equation}
\begin{equation}
Lu+\sigma_1u=0.
\tag{8.93}
\end{equation}

\begin{theorem}[Variational spectrum]
\label{gt-thm-8-37-variational-spectrum}
\dcref{ch8:8.37}
\group{ch08-generalized-solutions}
If the self-adjoint operator $L$ satisfies (8.5) and (8.6), then it has a
countably infinite discrete sequence of eigenvalues $\Sigma=\{\sigma_m\}$.
Writing $V_j$ for the eigenspace of $\sigma_j$, they are characterized by
\begin{equation}
\sigma_m=\inf\left\{J(u):u\in H\setminus\{0\},\
(u,v)_{L^2}=0\ \forall v\in V_1\oplus\cdots\oplus V_{m-1}\right\},
\tag{8.94}
\end{equation}
and their eigenfunctions form a complete system in $L^2(\Omega)$.
\end{theorem}
\begin{proof}
\sketch Lemma 8.4 bounds $J$ below. Normalize a minimizing sequence for
$\sigma_1=\inf_HJ$ in $L^2$. Coercivity bounds it in $H$, compact embedding
gives strong $L^2$ convergence, and the parallelogram identity plus Lemma 8.4
upgrades it to strong convergence in $H$. The Euler equation is
$Lu+\sigma_1u=0$. Repeat on the orthogonal complements of the preceding
eigenspaces. Compactness makes each eigenspace finite-dimensional, sends the
eigenvalues to infinity, and proves completeness.
\end{proof}

\begin{remark}[Eigenfunction regularity]
\label{gt-rem-8-eigenfunction-regularity}
\dcref{ch8:8.37}
\group{ch08-generalized-solutions}
The preceding regularity results imply that every eigenfunction belongs to
$L^\infty(\Omega)\cap C^\alpha(\Omega)$ for some $\alpha>0$, is H\"older
continuous up to the boundary under the hypotheses of Theorem 8.29, and is
smooth in $\Omega$ when the coefficients are smooth.
\end{remark}

\begin{theorem}[Simplicity of the first eigenvalue]
\label{gt-thm-8-38-simple-positive-first-eigenfunction}
\dcref{ch8:8.38}
\group{ch08-generalized-solutions}
Under the hypotheses of Theorem 8.37, the least eigenvalue is simple and has a
positive eigenfunction.
\end{theorem}
\begin{proof}
\sketch If $u$ minimizes the Rayleigh quotient, then $|u|$ does also. The Harnack
inequality makes $|u|$ positive almost everywhere. Every first eigenfunction
therefore has one sign; two linearly independent such eigenfunctions would yield
orthogonal first eigenfunctions, which is impossible. Hence the first eigenspace
is one-dimensional.
\end{proof}

\section*{Mixed Boundary Problems}

\begin{definition}[Generalized mixed boundary problem]
\label{gt-def-8-generalized-mixed-boundary-problem}
\dcref{ch8:8.96--8.97}
\group{ch08-generalized-solutions}
Let $\Gamma$ be a relatively open $C^1$ portion of $\partial\Omega$, let
$\nu=(\nu_1,\ldots,\nu_n)$ be the exterior unit normal on $\Gamma$, and put
$\Gamma_D=\partial\Omega\setminus\Gamma$. Consider
\begin{equation}
\begin{aligned}
Lu&=D_if^i+g &&\text{in }\Omega,\\
u&=\varphi_1 &&\text{on }\Gamma_D,\\
Nu:=a^{ij}\nu_iD_ju+b^i\nu_i u+\sigma u&=\varphi_2
&&\text{on }\Gamma.
\end{aligned}
\tag{8.96}
\end{equation}
Write $W_0^{1,2}(\Omega\cup\Gamma)$ for the closure of
$C_0^1(\Omega\cup\Gamma)$ in $W^{1,2}(\Omega)$. For
$\varphi_1\in W^{1,2}(\Omega)$ and $\sigma,\varphi_2\in L^2(\Gamma)$, a
generalized solution is a function $u\in W^{1,2}(\Omega)$ such that
$u-\varphi_1\in W_0^{1,2}(\Omega\cup\Gamma)$ and
\begin{equation}
\Qop(u,v)=\int_\Omega(f^iD_iv-gv)\,dx
+\int_\Gamma(\varphi_2-f^i\nu_i-\sigma u)v\,ds
\tag{8.97}
\end{equation}
for every $v\in W_0^{1,2}(\Omega\cup\Gamma)$.
\end{definition}

\begin{theorem}[Weak maximum principle for mixed boundary data]
\label{gt-thm-8-mixed-boundary-maximum-principle}
\dcref{ch8:8.98}
\group{ch08-generalized-solutions}
Assume (8.5), (8.6), and
\begin{equation}
\int_\Omega(dv-b^iD_iv)\,dx-\int_\Gamma\sigma v\,ds\leq0
\quad\forall v\geq0,\quad v\in C_0^1(\Omega\cup\Gamma).
\tag{8.98}
\end{equation}
If $u\in W^{1,2}(\Omega)$ satisfies
\[
\Lop(u,v)+\int_\Gamma\sigma uv\,ds\leq0
\]
for every nonnegative $v\in W_0^{1,2}(\Omega\cup\Gamma)$, then either
\[
\sup_\Omega u\leq\sup_{\Gamma_D}u^+
\]
or $u$ is a positive constant.
\end{theorem}
\begin{proof}
\sketch Repeat the truncation argument of Theorem 8.1 with test functions
vanishing on $\Gamma_D$. Condition (8.98) controls the interior and conormal
zeroth-order terms. The only exceptional case left by the support-measure
argument is a positive constant.
\end{proof}

\begin{corollary}[Uniqueness for the mixed problem]
\label{gt-cor-8-mixed-boundary-uniqueness}
\dcref{ch8:8.98}
\group{ch08-generalized-solutions}
Generalized solutions of (8.96) are unique if at least one of the following holds:
$\Gamma\neq\partial\Omega$, $\sigma+\nu_i b^i$ does not vanish identically on
$\Gamma$, or $L1\not\equiv0$. If all three conditions fail, any two solutions
differ by a constant. The associated existence statement follows from the same
Fredholm reduction as in Section 8.2, subject to the corresponding compatibility
condition when constants lie in the kernel.
\end{corollary}
\begin{proof}
\sketch Apply the mixed maximum principle to the difference of two solutions.
The three displayed conditions exclude its constant alternative. The Fredholm
argument uses $W_0^{1,2}(\Omega\cup\Gamma)$ in place of
$W_0^{1,2}(\Omega)$.
\end{proof}

\section*{Selected Further Consequences}

\begin{proposition}[Alternative hypotheses for the weak maximum principle]
\label{gt-prop-8-alternative-maximum-hypotheses}
\dcref{ch8:8.1}
\group{ch08-generalized-solutions}
In Theorem 8.1, when $u\leq0$ on $\partial\Omega$, condition (8.8) may be
replaced by either
\begin{equation}
\int_\Omega(dv+c^iD_iv)\,dx\leq0
\qquad\forall v\geq0,\quad v\in C_0^1(\Omega),
\tag{8.99}
\end{equation}
or the almost-everywhere block-matrix condition
\begin{equation}
\begin{bmatrix}
a&b\\
-c&-d
\end{bmatrix}\geq0.
\tag{8.100}
\end{equation}
\end{proposition}
\begin{proof}
\sketch For (8.99), use the adjoint form of the truncation computation. Under
(8.100), the combined principal, first-order, and zeroth-order terms in that
computation form a nonnegative quadratic expression.
\end{proof}

\begin{proposition}[Interior energy estimate]
\label{gt-prop-8-interior-energy-estimate}
\dcref{ch8:8.2}
\group{ch08-generalized-solutions}
If $u\in W^{1,2}(\Omega)$ solves $Lu=g+D_if^i$, where (8.5), (8.6) hold and
$g,f^i\in L^2(\Omega)$, then every $\Omega'\Subset\Omega$ satisfies
\begin{equation}
\|u\|_{W^{1,2}(\Omega')}
\leq C(\|u\|_2+\|f\|_2+\|g\|_2),
\tag{8.101}
\end{equation}
where
$C=C(n,\Lambda/\lambda,\nu,\dist(\Omega',\partial\Omega))$.
\end{proposition}
\begin{proof}
\sketch Test the weak equation against $\eta^2u$ for a cutoff supported a
controlled distance from the boundary. Ellipticity and Young's inequality absorb
the gradient terms; the cutoff derivative produces the stated distance dependence.
\end{proof}

\begin{proposition}[Symmetric-coefficient Harnack constant]
\label{gt-prop-8-symmetric-harnack-constant}
\dcref{ch8:8.3}
\group{ch08-generalized-solutions}
If $[a^{ij}]$ is symmetric, the constant in Theorem 8.20 may be chosen so that
\[
C\leq C_0^{\sqrt{\Lambda/\lambda+\nu R}},
\qquad C_0=C_0(n).
\]
\end{proposition}

\begin{proposition}[Sobolev extension and sharpness of the Poisson estimate]
\label{gt-prop-8-sobolev-poisson-sharpness}
\dcref{ch8:8.4}
\group{ch08-generalized-solutions}
\uses{gt-ch3-poisson}
The conclusion of Theorem~\ref{gt-ch3-poisson} remains valid for weak
$W^{1,2}(\Omega)$ solutions of Poisson's equation with bounded right-hand side.
In
\[
\Omega=\{(x,y)\in\mathbb R^2:|x|+|y|<1\},
\]
the function
\[
u(x,y)=|xy|\log(|x|+|y|)
\]
shows that the logarithmic modulus in that theorem is sharp.
\end{proposition}
\begin{proof}
\sketch Regularize the weak solution on compact subdomains, apply
Theorem~\ref{gt-ch3-poisson} to the regularizations, and pass to the limit using
Theorem 8.8. Direct differentiation of the displayed planar example gives the
borderline logarithmic modulus.
\end{proof}

\begin{lemma}[Planar circle-oscillation estimate]
\label{gt-lem-8-planar-circle-oscillation}
\dcref{ch8:8.5(a)}
\group{ch08-generalized-solutions}
Let $u\in C^1(B_R(0))$, where $B_R(0)\subset\mathbb R^2$, and set
\[
\omega(r)=\osc_{\partial B_r}u,
\qquad
D(r)=\int_{B_r}|Du|^2\,dx.
\]
If $\omega$ is nondecreasing, then, for $0<r<R$,
\[
\omega(r)\leq\sqrt{\frac{\pi D(R)}{\log(R/r)}}.
\]
\end{lemma}
\begin{proof}
\sketch Bound the oscillation on almost every circle by the tangential
Dirichlet energy and Cauchy--Schwarz. Integrate the resulting lower bound for
$D'(s)$ over $r<s<R$ and use monotonicity of $\omega$.
\end{proof}

\begin{lemma}[Capacity controlled by measure]
\label{gt-lem-8-capacity-measure}
\dcref{ch8:8.7}
\group{ch08-generalized-solutions}
For every bounded measurable $\Omega\subset\mathbb R^n$,
\[
|\Omega|^{1-2/n}\leq\gamma(n)\,\capop\Omega,
\]
where $\gamma(n)$ is the constant in the Sobolev inequality (7.26) for $p=2$.
\end{lemma}
\begin{proof}
\sketch Apply the Sobolev inequality to each compactly supported function that
is at least one on $\Omega$, then take the infimum of its Dirichlet energy.
\end{proof}

\begin{theorem}[Finite-measure unbounded domains]
\label{gt-thm-8-finite-measure-unbounded-dirichlet}
\dcref{ch8:8.9}
\group{ch08-generalized-solutions}
The existence and uniqueness conclusion of Theorem 8.3 remains valid when
$\Omega$ is unbounded but has finite measure.
\end{theorem}
\begin{proof}
\sketch Theorem 8.16 supplies the a priori control needed to replace the compact
embedding used in the bounded-domain Fredholm proof. Exhaust $\Omega$ by
bounded subdomains, solve there, and pass to a weak limit; the maximum principle
gives uniqueness.
\end{proof}
